\subsection{Divisibility of polynomials in one indeterminate over a field}

The polynomial ring K{[}X{]} in one indeterminate over a commutative field K is a principal ideal domain (IV, p.~11, Prop. 11). Its field of fractions is the field K(X) of rational functions in X with coefficients from K. The ring K{[}X{]} contains the subring of polynomials of degree 0, that is the field of constants, which is identified with K; the elements of K* are invertible in K, and hence in K{[}X{]}; conversely the formula \(\deg(uv) = \deg(u) + \deg(v)\) shows that every invertible polynomial has degree 0; the group U of invertible elements of K{[}X{]} is thus precisely K*. Thus two associate polynomials differ only by a nonzero constant factor; in particular every class of associate polynomials contains a unique monic polynomial. The subgroup of the multiplicative group K(X)* generated by the monic polynomials thus contains a unique element from each class of associate rational functions, and consequently is isomorphic to the group

\[
\mathcal {P} ^ {*} = \mathrm{K} (\mathrm{X}) ^ {*} / \mathrm{U}
\]

of principal fractional ideals of \(\mathrm{K}(\mathrm{X})\) . In particular, whenever gcd or lcm in the field \(\mathrm{K}(\mathrm{X})\) (with respect to the ring \(\mathrm{K}[\mathrm{X}]\) ) is mentioned, it will usually be understood that these are quotients of monic polynomials (or 0); this convention allows us to speak of the gcd or the lcm of a family of rational functions.

The irreducible elements of K{[}X{]} are precisely the irreducible polynomials in the usual sense (IV, p.~13, Def. 2), and the set of monic irreducible polynomials is a system of representatives of irreducible elements of K{[}X{]}.

A first degree polynomial is always irreducible. If K is an algebraically closed field then the converse is true (V, p.~19, Prop. 1); thus in this case every polynomial \(p(\mathbf{X})\) of degree \(n\) in \(\mathbf{K}[\mathbf{X}]\) can be written uniquely (up to the order of the factors) in the form

\[
p (\mathbf {X}) = c \left(\mathbf {X} - a _ {1}\right) \left(\mathbf {X} - a _ {2}\right) \dots \left(\mathbf {X} - a _ {n}\right)
\]

where \(c\) and the \(a_{i}\) are elements of \(\mathbf{K}\) .

PROPOSITION 6. --- For any field K, the set of monic irreducible polynomials in K{[}X{]} is infinite.

Indeed, given an arbitrary nonempty finite family \((p_{i})\) ( \(1 \leqslant i \leqslant n\) ) of distinct monic irreducible polynomials, the polynomial \(\left(\prod_{i=1}^{n} p_{i}\right) + 1\) is not invertible, and any monic irreducible factor q of this polynomial is necessarily distinct from all the \(p_{i}\) , otherwise it would divide 1.

